\section{Appendix: Scope and Qualifications}
\label{sec:scope}

{\small
\textbf{Edition and formulary.} The study treats the Mass of the
Twenty-eighth Sunday in Ordinary Time, Year~A, as the \work{Roman Missal},
Third Edition, and the \work{Lectionary for Mass} give it for the dioceses of
the United States, in the national calendar alone, for the occurrence of
11~October 2026. Identity, occurrence and branches are set out in the
preceding appendix and recorded in \texttt{instance/manifest.md} and
\texttt{research/context.md}.

\textbf{Text control.} The Missal formulary was read in the Latin of the
\work{Missale Romanum}, editio typica tertia (2002), in a digitally typeset
reproduction, and compared with the list of changes for the 2008 reprint in
\work{Notitiae} 44 (2008), which lists none for this week, and with the
\work{Antiphonary} excerpted from the English Missal (2010). No page of the
2008 reprint and no United States altar book was inspected, so the approved
English orations were not read in any witness, and no chant book was
collated. The Latin source of the Entrance Antiphon's closing words is not
established. The Lectionary's boundaries rest on three witnesses that agree in
every figure: the national calendar for 2026, the bishops' conference's dated
page of readings, and no.~142 of the Latin \work{Ordo lectionum Missae} (1981);
the printed United States Lectionary was not inspected. Four questions about
the text remain open: whether the printed Lectionary gives the psalm response
as the dated page does; whether the United States Lectionary admits another
Alleluia verse from a common set on these Sundays; the Greek of Isaiah 25:6--10
and of Philippians 4:19, which was not examined; and manuscript evidence for
Matthew 22:1--14 beyond what the SBL apparatus reports. No manuscript was
examined.

\textbf{English.} All Scripture is printed in the Douay--Rheims version as
revised by Bishop Challoner, a public-domain study translation of the
Clementine Vulgate. It is not the text proclaimed in the celebration; the
controlling English is that of the \work{Lectionary for Mass for Use in the
Dioceses of the United States of America}, and for the antiphons and orations
that of the \work{Roman Missal}, Third Edition, neither of which is reproduced.
The orations and the antiphons are identified by Latin incipit and described;
the descriptions are not translations and may not be quoted as the texts'
words. Where the Lectionary appoints part of a verse, the rest is marked
beside it. English glosses of Latin passages are the editor's renderings of
the Latin printed beside them.

\textbf{Reception.} Each appointed Scripture passage has at least one ancient
and one medieval or later witness read at an exact locus; the witnesses are
those named in the References, and the search behind them is bounded and
recorded in \texttt{research/scope.md}. For the second Communion antiphon the
later witness is Aquinas's doctrinal use of the verse. Five passages rest on
one ancient commentator each: the Entrance Antiphon's psalm and both Communion
antiphons on Augustine, and the second reading and the source of the Alleluia
verse on Chrysostom. Chrysostom, Augustine's sermons, his
expositions of Psalms 22, 33 and 129 and his homilies on the First Epistle of
John, the \work{Catechetical Lectures} of Cyril of Jerusalem and Irenaeus were
read in the English of the Nicene and Ante-Nicene Fathers, Irenaeus in a
machine-read text of the printed volume; Aquinas's \work{Summa} I in the
English Dominican translation; Bellarmine in O'Sullivan's abridged English of
1866, the Latin not compared; Schuster in a text layer of the 1927 English and
Fromage in the English of 1909, on a page image. Jerome on Matthew and Isaiah, Hilary on Matthew,
and Aquinas on Matthew, Isaiah, Ephesians and Philippians were read on page
images of the printings named; Gregory, Aquinas's commentary on the Psalms and
\work{Summa} I-II in electronic Latin texts. No witness was collated with a
modern critical edition. Not opened, among others: Origen and Rabanus Maurus
on Matthew; Eusebius, Cyril of Alexandria and Theodoret on Isaiah; Cassiodorus
on the psalms; Basil on Psalm 33; Marius Victorinus and Ambrosiaster on
Philippians; Jerome on Ephesians; Bede on the First Epistle of John. Nothing is
said here about who first proposed any reading, about any reading being the
Fathers' common view, or about the order in time or the mutual dependence of
the witnesses. No commentary on the Prayer over the Offerings or the Prayer
after Communion was found. Fromage's and Schuster's comments on the Collect,
and Schuster's on the Gospel, come from their commentaries on Masses that
differ from this formulary in their other elements, and each is cited only
for the element it shares with this one. The \work{Mystagogical Catecheses}
are cited as transmitted under Cyril's name: their editor reports manuscripts
that give them to John of Jerusalem, his successor, and they support no claim
here on their own.

\textbf{Relations among the elements.} Four relations are stated by the
liturgical books: the first reading is chosen in relation to the Gospel
(\work{Praenotanda} 106), the psalm responds to the first reading
(\work{General Instruction} 61), the acclamation greets the Gospel
(\work{General Instruction} 62), and the second reading and the Gospel each
belong to a semi-continuous course. The Missal's antiphons and orations serve
all three years. The recurrence of the words for calling, feasting and hunger,
riches and glory, clothing and seeing across the texts is an observation about
the texts. Among the witnesses cited, Aquinas joins Isaiah 25:6 to Matthew
22:4 and to Psalm 22:5, and reads Psalm 22:6 as the subsequent grace whose
pair of terms the Collect uses, without citing the Collect; Augustine joins
Psalm 33:11 to John 6:51. Every other connection drawn between one element and
another is the editor's synthesis, grounded in a witness's statement about his
own passage where one is cited and otherwise an editorial proposal resting on
the wording of the texts. No witness cited comments on this Mass.

\textbf{The Council's declaration.} \work{Nostra aetate} 4 was read in the
English and Latin of the Holy See's website. It is cited as the Church's
teaching on how the identification of the first invited with a people is
presented and applied, not as a judgment on the exegesis of any Father.

\textbf{Chronology.} The Date cells print the shared chronology corpus's
answers with their relations and alternatives unchanged. The first reading's
answer dates the prophet's ministry in one author's account, not the oracle.
The Psalter's common composition bound is not a date for any of the three
psalms; it was read in the current public text of the introduction it cites,
not in the registered state. The superscription answer of Psalm 33 dates the
episode its title names and rests on Haydock's notes, whose pages were not
read here. The Gospel's event answer dates the telling of the parable in one
author's chronology; its composition answers, and those of the epistles, are
dates of the books. Writing sites, first readers and narrative settings are
kept apart; the source checks and their limits are recorded in
\texttt{research/scope.md}.

\textbf{Rights.} The Douay--Rheims, the Clementine Vulgate, the King James
Version, the Nicene and Post-Nicene and Ante-Nicene Fathers, O'Sullivan's
Bellarmine, the English Dominican \work{Summa}, the volumes of Migne, the
Venice, Parma and Li\`ege printings of Aquinas, Fromage's and Schuster's
English, Wilson's and Feltoe's editions and the \work{Catholic Encyclopedia}
are in the public domain. Gregory's Latin is quoted in short passages from a
licensed electronic edition, and the Corpus Thomisticum texts of Aquinas's
commentary on the Psalms and of \work{Summa} I-II in short phrases. The Latin
and English of the Missal, the \work{Ordo lectionum Missae}, the Lectionary,
the \work{General Instruction}, the national calendar, the \work{Nova
Vulgata}, \work{Nostra aetate} and the introduction to the Psalms of the New
American Bible are protected; they are cited, and only incipits, the three
titles of the \work{Ordo}, short phrases of the Council's declaration and a few
words needed by the argument are quoted.

\textbf{Records.} The composition of the target is audited in
\texttt{propers/verified.md}; the reception search, loci, disagreements and
negative results in \texttt{research/scope.md}; and the reasoning behind the
three interpretations, with the class of every cross-element link, in
\texttt{research/interpretations.md}. The research records received an
independent review before the study was written.
\par}

\phantomsection
\section*{References}
% A starred heading sets no mark, so without this the Appendix above would go
% on heading the pages of the bibliography. The heading stays starred because
% check-content-preflight finds the References section by that exact form.
\runninghead{References}
\addcontentsline{toc}{section}{References}
\label{sec:references}

{\small
\subsection*{Liturgical books, Bibles and texts}
\begin{itemize}[itemsep=0pt,topsep=.25em]
\item \work{Missale Romanum}, editio typica tertia (Typis Vaticanis, 2002), \latin{Dominica XXVIII ``per annum''}; the rubrics of \latin{Tempus per annum} 1--6; the \latin{Normae universales de anno liturgico et de calendario} 59--60, as printed in that Missal.
\item \work{Variationes et additiones in Missali Romano}, \work{Notitiae} 44 (2008), pp.~371--372.
\item \work{The Roman Missal}, Third Edition, for Use in the Dioceses of the United States of America (2011), as represented by the \work{Antiphonary: Excerpted from the Roman Missal} (International Commission on English in the Liturgy, 2010), p.~80.
\item \work{Ordo lectionum Missae}, editio typica altera (Libreria Editrice Vaticana, 1981): \work{Praenotanda} 80, 105--107; \latin{Tabella II}; no.~142, \latin{Dominica vigesima octava}, Year~A.
\item \work{Lectionary for Mass for Use in the Dioceses of the United States of America}, second typical edition, vol.~I, no.~142, as witnessed by the United States Conference of Catholic Bishops, \work{Liturgical Calendar for the Dioceses of the United States of America} 2026, pp.~5 and 40, and its page of readings for 11~October 2026, \url{https://bible.usccb.org/bible/readings/101126.cfm}.
\item \work{General Instruction of the Roman Missal}, United States edition of 2011 with the emendations of 2021, nos.~48, 51, 53, 61--64, 68, 74, 87, 363--365.
\item The Holy Bible, Douay--Rheims version, revised by Bishop Richard Challoner (Project Gutenberg text): 1~Kgs 21:10--15; Ps 16:15; 18:6; 22; 33:1, 9--12; 58:11; 62:6; 71:2; 129; 138:16; Is 1:1; 24:23--25:12; Jer 9:23; Mt 21:23--22:15; Lk 14:18--20; Rom 11:29; 13:14; 1~Cor 11:26; 13:12; 15:53--55; 2~Cor 5:4; Gal 3:27; Eph 1:1, 15--19; 3:1; Phil 1:1, 7, 13; 4:10--23; Col 3:14; 1~Tim 1:5; 1~Pet 3:21; 2~Pet 1:3--4; 1~Jn 3:1--3.
\item \work{Biblia Sacra Vulgatae editionis} (Clementine Vulgate), eBible.org text: Ps 22:5--6; 33:11; 71:2; 129:3--4; Is 25:6--10; Mt 22:2--14; Eph 1:17--18; Phil 4:12--13, 19; 1~Pet 3:21; 2~Pet 1:4; 1~Jn 3:2.
\item \work{Nova Vulgata Bibliorum Sacrorum editio}, as published by the Holy See at vatican.va: Ps 34 (33):11; 130 (129):4; Is 25:6--9; Phil 4:19; 1~Jn 3:2.
\item The Holy Bible, King James Version (eBible.org text): Ps 34:10.
\item \work{The Greek New Testament: SBL Edition}, ed. M.~W. Holmes, version~1.2, text and apparatus of Matthew at 22:1--14.
\end{itemize}

\subsection*{Fathers, Doctors and saints}
\begin{itemize}[itemsep=0pt,topsep=.25em]
\item St Irenaeus of Lyons, \work{Against Heresies} IV.36.5--6, Ante-Nicene Fathers, vol.~1 (Buffalo, 1887), pp.~516--517, in the machine-read text of the printed volume.
\item St Hilary of Poitiers, \work{Commentarius in Matthaeum} 22.3--7 (Migne, PL 9, Paris, 1844, cols.~1042--1044).
\item St John Chrysostom, \work{Homilies on Matthew} 69.1--4, Nicene and Post-Nicene Fathers, 1st ser., vol.~10; \work{Homilies on Ephesians} 3 and \work{Homilies on Philippians} 15, Nicene and Post-Nicene Fathers, 1st ser., vol.~13 (both in the electronic text of the Christian Classics Ethereal Library).
\item St Cyril of Jerusalem, \work{Procatechesis} 3; and the \work{Mystagogical Catecheses} transmitted under his name, 3.4, 3.7, 4.3, 4.7; Nicene and Post-Nicene Fathers, 2nd ser., vol.~7.
\item St Jerome, \work{Commentarii in Matthaeum} III, on 22:1--14 (Migne, PL 26, Paris, 1845, cols.~159--161); \work{Commentarii in Isaiam} VIII, on 25:1--12 (Migne, PL 24, Paris, 1865, cols.~300--301).
\item St Augustine, \work{Sermo} 90.1, 4--6 and 95.5--7 (Sermons 40 and 45 on the New Testament Lessons), Nicene and Post-Nicene Fathers, 1st ser., vol.~6.
\item St Augustine, \work{Enarrationes in Psalmos} 22.1--6, 33 (sections 13--14 of the English text) and 129.1--3, Nicene and Post-Nicene Fathers, 1st ser., vol.~8.
\item St Augustine, \work{In epistolam Iohannis ad Parthos tractatus} 4.5--7 (Homilies on the First Epistle of John), Nicene and Post-Nicene Fathers, 1st ser., vol.~7.
\item St Gregory the Great, \work{Homiliae in Evangelia} 38.1, 3--7, 9--12, 14, 16, Latin (Wikisource text).
\item St Thomas Aquinas, \work{Super Evangelium Matthaei} c.~22, in \work{Opera}, editio altera Veneta, vol.~3 (Venice: Bettinelli, 1745), pp.~279--282.
\item St Thomas Aquinas, \work{Expositio super Isaiam ad litteram} c.~25, in \work{Opera omnia}, vol.~14 (Parma, 1863), pp.~501--502.
\item St Thomas Aquinas, \work{Super epistolam ad Ephesios} c.~1, lect.~6, and \work{Super epistolam ad Philippenses} c.~4, lect.~2, in \work{In omnes S.~Pauli Apostoli epistolas commentaria}, vol.~2 (Li\`ege: Dessain, 1857), pp.~272 and 401--403.
\item St Thomas Aquinas, \work{Postilla super Psalmos} 22 (Corpus Thomisticum text).
\item St Thomas Aquinas, \work{Summa theologiae} I, q.~12, aa.~1, 6 (English Dominican translation, Project Gutenberg text); I-II, q.~110, a.~3; q.~111, a.~3; q.~112, a.~1 (Corpus Thomisticum text).
\item St Robert Bellarmine, \work{A Commentary on the Book of Psalms}, trans. J. O'Sullivan (1866, abridged), in the eCatholic2000 transcription: on Ps~22:1--6, Ps~33:10--11 and Ps~129:1--6.
\end{itemize}

\subsection*{Liturgical history and commentary}
\begin{itemize}[itemsep=0pt,topsep=.25em]
\item Lucien Fromage, continuation of Prosper Gu\'eranger, \work{The Liturgical Year}, vol.~11, Time after Pentecost II, trans. Laurence Shepherd, 2nd ed. (1909), p.~357.
\item Blessed Ildefonso Schuster, \work{The Sacramentary (Liber Sacramentorum)}, vol.~3 (London: Burns Oates \& Washbourne, 1927), pp.~142 and 173.
\item H.~A. Wilson, ed., \work{The Gregorian Sacramentary under Charles the Great} (London, 1915), pp.~61, 135, 167 and 174.
\item C.~L. Feltoe, ed., \work{Sacramentarium Leonianum} (Cambridge, 1896), in a machine-read text, for the Prayer after Communion.
\end{itemize}

\subsection*{Magisterium}
\begin{itemize}[itemsep=0pt,topsep=.25em]
\item Second Vatican Council, Declaration on the Relation of the Church to Non-Christian Religions \work{Nostra aetate} (28~October 1965), no.~4, English and Latin texts of the Holy See's website.
\end{itemize}

\subsection*{Chronology sources}
\begin{itemize}[itemsep=0pt,topsep=.25em]
\item \work{The Catholic Encyclopedia} (New York, 1908--1912): Corbett, ``David'' (vol.~4, 1908); ``Epistle to the Ephesians'' (vol.~5, 1909); Souvay, ``Isaias'', Maas, ``Jesus Christ'', and Drum, ``Epistles of Saint John'' (vol.~8, 1910); Jacquier, ``Gospel of St.~Matthew'' (vol.~10, 1911); Prat, ``St.~Paul'' (vol.~11, 1911); Vander Heeren, ``Epistle to the Philippians'' (vol.~12, 1911); Durand, ``The New Testament'' (vol.~14, 1912); read in the New Advent text.
\item United States Conference of Catholic Bishops, \work{New American Bible, Revised Edition}, introduction to the Psalms, \url{https://bible.usccb.org/bible/psalms/0}, read in its public delivery of 8~October 2026 (summarized; protected).
\item George Leo Haydock, notes on 1~Kings 21 and Psalm~70, in the Douay--Rheims Bible with Haydock's commentary, as the chronology corpus records them for the A.M.\ figure.
\end{itemize}
\par}
